%%%%%%%% SD4H @ ICML 2026 — Camera-ready Path B (closure-canonical reframe) %%%%%%%%%%%%%%%%
% Path B: reconcile Results to the 2026-06-01 Pipeline 2 closure.
%   - Single canonical model = 2-Component cortical opponent-axis (2-DOF).
%   - Machado / R+C presented as candidate mechanisms that FAIL structural
%     adequacy; R+C rejected (boundary saturation).
%   - Selection = held-out test_loss_median + structural adequacy (NOT perm-p).
%   - HC specificity = descriptive only; physiological (beta_s,beta_c)
%     interpretation removed (noise floor; 0/3 dual-pass null).
%   - Divergence thesis anchored on sub-08 (identified beta_c sign); both
%     mechanisms feasible (8/8), corrections oppositely directed (cos ~ -0.4,
%     sign-agree 1/8). sub-09 NOT a divergence case: its beta_c sign is
%     non-identified, so its retinal-vs-cortical cosine flips with the sign.
%     "4/8 collapse / 360->96 arc" REMOVED (superseded c3_relabel; both 8/8).
%   - Base = SD4H_draft_v6.1.tex; preamble/title/packages/citation keys preserved.
%
% v6.1: Concrete-first framing pass over v6.
%   - Abstract/intro open with "such as CVD" instead of pure abstraction.
%   - "distortion field" terminology removed throughout
%     ("the fitted model" at L92, "low-dimensional distortion" at L93).
%   - Representational-geometry sentence split: general definition,
%     then CVD-specific application.
%   - Results: significance vs. permutation null stated once at the
%     §3.1 opener and once in Fig 2 caption; per-subject blurbs carry
%     numeric ρ / p_perm only.

\documentclass{article}

\usepackage{microtype}
\usepackage{graphicx}
\usepackage{subcaption}
\usepackage{booktabs}
\usepackage{hyperref}

\newcommand{\theHalgorithm}{\arabic{algorithm}}

\usepackage{icml2026}

\usepackage{amsmath}
\usepackage{amssymb}
\usepackage{mathtools}
\usepackage{xcolor}
\usepackage{multirow}

\usepackage[capitalize,noabbrev]{cleveref}

% Inline TODO marker for data-dependent values to be filled before submission
\newcommand{\todo}[1]{\textcolor{red}{[TODO: #1]}}

% Placeholder command for figures not yet created
\newcommand{\figplaceholder}[2]{%
  \begin{center}
    \fbox{\parbox{#1}{\centering\vspace{1.5em}\textcolor{gray}{\textbf{[FIGURE PLACEHOLDER]}\\#2}\vspace{1.5em}}}
  \end{center}
}

\icmltitlerunning{Inferring Individualized Color-Vision Distortions from fMRI Hue-Representation Geometry}

\begin{document}

\twocolumn[
  \icmltitle{Inferring Individualized Color-Vision Distortions from \\
    fMRI Hue-Representation Geometry}

  \icmlsetsymbol{equal}{*}

  \begin{icmlauthorlist}
    \icmlauthor{Anonymous Authors}{anon}
  \end{icmlauthorlist}

  \icmlaffiliation{anon}{Anonymous Institution}

  \icmlcorrespondingauthor{Anonymous Author}{anonymous@example.com}

  \icmlkeywords{structured representation, individualized phenotype inference, representational geometry, color vision deficiency, fMRI}

  \vskip 0.3in
]

\printAffiliationsAndNotice{}

%==============================================================================
\begin{abstract}
%==============================================================================
Many health-related data carry structured distortions that can act as phenotypes, yet interpreting such signatures and translating them into intervention remains challenging.
Color vision deficiency (CVD), for example, distorts pairwise relations between neural hue representations while categorical recognition remains relatively preserved.
Here we formulate fMRI-measured hue responses as structured neural representations and infer low-dimensional distortions that explain their geometry.
We quantify this distortion using the Leave-One-Color-Out (LOCO) interpolation vulnerability profile, which measures how well each held-out hue can be interpolated from the remaining hue manifold.
Applied to Ishihara-confirmed CVD participants across visual areas (V1--hV4), the framework recovers each subject's distortion with a two-degree-of-freedom cortical opponent-axis model, selected among candidate retinal and retinal--cortical mechanisms by held-out predictive loss and structural adequacy.
Candidate mechanisms agree at the level of which individuals are distorted, yet this detection-level agreement cannot adjudicate the correction: the retinal and cortical accounts both yield feasible corrections, but in nearly opposite directions, so only structural adequacy and a planned behavioral test can select between them.
These findings show that detection-level agreement is insufficient for choosing an intervention and demonstrate that structured distortions in high-dimensional biological data can be resolved into low-dimensional, subject-specific models for personalized intervention.
\end{abstract}

%==============================================================================
\section{Introduction}
\label{sec:intro}
%==============================================================================

Many health-related data carry distorted structure that can act as a phenotype—systematic, repeatable patterns in how stimuli relate within a high-dimensional biological response space.
Color vision deficiency (CVD), for example, produces distorted pairwise relations between neural hue representations even while categorical recognition remains relatively preserved.
Other instances of such structured-phenotype data include geometric deformations of cardiac anatomy \citep{biffi2020} and directional shifts in gene-expression latent spaces under drug perturbation \citep{lotfollahi2019}.
Yet such phenotypes are challenging to interpret and, especially, to translate into individually-targeted intervention.
When such structural distortion can be parameterized, it can be inverted for correction—a route absent from current generic correction approaches.

In a geometric framework, responses to stimuli can be summarized as a \emph{representational geometry}, where pairwise distances capture what a neural population distinguishes \citep{kriegeskorte2008}.
In CVD, this geometry becomes the natural inference target: the pairwise relationships among hues can be systematically distorted even when categorical decoding remains preserved.
Despite this perspective, structured distortions have rarely been used to derive individualized perceptual corrections; the closest analogue---a cochlear simulator that inverts a hearing-impairment deformation model \citep{irino2023}---targets simulation of sensory loss, not its correction.

CVD arises from altered spectral sensitivity of retinal cone photoreceptors, affecting approximately 8\% of males \citep{machado2009}, and produces structured distortions in cortical color representations that vary across individuals.
Current assistive approaches, including commercial spectral-notch filters, apply a fixed, generic correction; recent evidence indicates that such filters alter color \emph{appearance} but have minimal effect on color \emph{discrimination} at threshold \citep{somers2024}, suggesting that effective correction may require individually-tailored strategies.

While categorical color decoding can remain preserved---leave-one-run-out decoding shows no HC--CVD difference ($p = 0.668$)---continuous hue relationships may be systematically altered; classification accuracy is insensitive to this relational structure.
CVD is also heterogeneous across individuals; even within a single diagnostic category, cortical compensation can vary substantially \citep{tregillus2021}, motivating individual-level rather than group-level inference.
In CVD, the pairwise geometry of stimulus-evoked fMRI responses is itself distorted in low-dimensional, parameterizable ways; we model this distortion as a mapping that can be fit and inverted.

Here we present a framework that quantifies this structured distortion through the interpolation vulnerability profile—how well each hue can be predicted from its neighbors—and fits low-dimensional simulator parameters that reproduce the observed profile.
The fitted model can then be used to generate candidate corrections through pre-image search.
The unit of inference is each individual's distortion model, not the population mean.
Our contributions are:
\begin{enumerate}
    \item We formulate distorted neural geometry as a structured-distortion inference problem and quantify it through the interpolation vulnerability profile, which captures subject-specific deficits in pairwise hue relations beyond classification accuracy.
    \item We infer a low-dimensional (2-DOF) cortical opponent-axis distortion model from high-dimensional neural representations, selecting it among candidate retinal and retinal--cortical mechanisms by held-out predictive loss and structural adequacy rather than by in-sample significance.
    \item We show that candidate mechanisms agree on which individuals are distorted yet prescribe divergent corrections, so that detection-level agreement cannot determine the correction; this gap defines a concrete, falsifiable behavioral test as the framework's next step.
\end{enumerate}

%==============================================================================
\section{Methods}
\label{sec:methods}
%==============================================================================

\begin{figure}[t]
  \vskip 0.1in
  \centering
  \includegraphics[width=\columnwidth]{figures/fig1_panels_bcd.pdf}
  \caption{Structured distortion and the inference pipeline.
  \textbf{Top}: CVD distorts the continuous geometry of hue representations in hV4---some neighboring hues collapse while others expand---even though classification accuracy is preserved.
  \textbf{Bottom}: The inference pipeline quantifies the structured distortion via the LOCO vulnerability profile, fits parsimonious distortion parameters (1--2 DOF), and assesses correction feasibility via pre-image inversion.}
  \label{fig:pipeline}
  \vskip -0.1in
\end{figure}

\subsection{Data Acquisition and Structured Representation}

Ten participants (7 healthy controls [HC], 3 females, age $22.7 \pm 2.5$ years; 3 color vision deficient [CVD], all male: two deutan and one protan) were assessed using the 14-plate Ishihara test \citep{ishihara1917}.
All three CVD participants are Ishihara-confirmed; one deutan case (Sub-10) is null at both the LOCO and SRM representational levels used in this work and is therefore reserved as a future-work probe of alternative representational dimensions.
Given the small CVD sample, all CVD results are interpreted as individual case demonstrations, not population-level effects \citep{crawford1998}.

Stimuli consisted of eight colors equally spaced from 0$^\circ$ to 315$^\circ$ in CIE~$L^*a^*b^*$ space ($L^* = 75$, chroma distance 40).
Participants performed a modified Rapid Serial Visual Presentation (RSVP) task \citep{brouwer2009} across six runs, viewing 1.5~s color patches.
Functional images (Siemens 3T; TR\,=\,1.5~s; 2~mm isotropic; 24 oblique slices) were preprocessed with fMRIPrep and projected to retinotopically-defined ROIs (V1, V2, hV4) via the Wang probabilistic atlas at 50\% threshold \citep{wang2015}.
Response amplitudes were estimated using a two-stage FIR-then-GLM procedure \citep{brouwer2009} and Procrustes-aligned across runs.

We modeled hue-selective neural populations with half-wave rectified and squared sinusoidal basis functions \citep{brouwer2009}: $r_k(\theta) = \max(0, \cos(\theta - \theta_k))^2$ for $K$ channels.
The weight matrix $W$ mapping channels to voxel responses was estimated via ridge regression with GCV regularization \citep{golub1979}.
This encoding model transforms raw voxel data into a structured representation in which the relational geometry across hues is explicitly parameterized.

\subsection{Structured Distortion Quantification: Interpolation Vulnerability}

To capture distortions in relational geometry, we extract an interpolation vulnerability profile that reveals which hues are geometrically fragile in a participant's neural representation.
In leave-one-color-out (LOCO) cross-validation, the encoding model $W$ is trained on 7 of 8 hues and evaluated on the held-out hue; the 8 prediction errors form a vulnerability vector $\mathbf{v} \in \mathbb{R}^8$, capturing structural deficits invisible to categorical classification.

This profile is specific to continuous geometry: leave-one-\emph{run}-out (LORO) classification shows no HC--CVD difference, while LOCO interpolation error is significantly elevated in CVD.
Among the ROIs tested, the effect was most robust in hV4 ($p = 0.017$; $K{=}3$), consistent with hV4's established role in perceptual color representation \citep{brouwer2009, bannert2018}.
We therefore focus distortion inference on hV4, with V1 analyses reported as cross-ROI validation in \cref{app:supporting}.

We additionally quantify pairwise structural distortion via representational distance matrices (RDMs).
The $8 \times 8$ crossnobis distance matrix is computed in voxel space for each participant, and the difference $\Delta\text{RDM} = \text{RDM}_{\text{CVD}} - \overline{\text{RDM}}_{\text{HC}}$ measures how a CVD individual's pairwise geometry deviates from the HC mean.
While $\mathbf{v}$ summarizes per-color vulnerability (8-dim), $\Delta$RDM captures pairwise distortion ($\binom{8}{2} = 28$-dim)---used as a fitting regularizer (\cref{eq:loss}) and cross-criterion validation target (\cref{sec:results}).

\subsection{Inverse Inference of Low-Dimensional Distortion Fields}

We infer low-dimensional distortion parameters from the observed profile $\mathbf{v}$ by evaluating three candidate mechanisms that parameterize how a stimulus hue angle $\theta$ maps to a distorted angle $\theta'$ at distinct processing levels (\cref{fig:pipeline}, bottom; notation in \cref{app:notation}): a retinal cone shift (Machado), a retinal shift under cortical compensation (R+C), and a cortical opponent-axis reorganization (2-Component).
These are competing hypotheses about the locus of distortion, not a nested escalation; we ask which mechanism is structurally adequate to reproduce the observed geometry and generalizes to held-out data (\cref{sec:methods_select}).

\paragraph{Machado cone shift (1 DOF).}
A retinal model \citep{machado2009}: the anomalous cone's spectral peak is shifted by $\Delta\lambda \in [0, 20]$~nm, predicting hue distortion through altered L--M cone opponency.

\paragraph{Retinal + cortical compensation (R+C; 1 DOF).}
Scales the Machado retinal shift by a cortical compensation factor:
\begin{equation}
\delta\theta_{\text{RC}}(c) = (2 - g) \cdot \delta\theta_{\text{Machado}}(c; \Delta\lambda)
\label{eq:rc}
\end{equation}
where $\delta\theta_{\text{Machado}}$ is the per-color hue shift induced by cone shift $\Delta\lambda$ and $g$ scales the cortical response: $g{=}1$ is no compensation, $g{=}2$ is exact compensation, and $g{>}2$ is overcompensation---motivated by evidence for progressive cortical amplification in anomalous trichromats \citep{tregillus2021}.
Because R+C rescales the retinal shift along a single axis, it carries no confusion-axis degree of freedom; we show below (\cref{sec:methods_select}, \cref{app:supporting}) that this structural limitation causes $g$ to saturate at the grid boundary, rejecting R+C as structurally inadequate.

\paragraph{Cortical opponent-axis model (2 DOF).}
The cortical mechanism models distortion as angular dilation along two opponent axes:
\begin{equation}
\theta'(c) = \theta(c) + \underbrace{\beta_s \cos(\theta(c) {-} 90^\circ)}_{\text{S-cone axis}} + \underbrace{\beta_c \cos(\theta(c) {-} \theta_{\text{conf}})}_{\text{confusion axis}}
\label{eq:2comp}
\end{equation}
where $\beta_s$ captures S-cone axis expansion, motivated by behavioral evidence for hue-angle reorganization in CVD \citep{emery2021}, and $\theta_{\text{conf}} = 16^\circ$ (protan) or $150^\circ$ (deutan) is the type-specific confusion axis.

\subsection{Model Selection and Specificity}
\label{sec:methods_select}

We select optimal parameters by forward-simulating each model's distortion through the HC encoding model and comparing the resulting profile to the observed $\mathbf{v}$.
The composite fitting loss balances per-color functional accuracy with pairwise structural coherence:
\begin{equation}
L = \underbrace{L_{\text{vuln}} + 0.5 \, L_{\text{rank}}}_{\text{per-color (8-dim)}} + \underbrace{0.2 \, L_{\text{rdm}}}_{\text{pairwise (28-dim)}} + 0.1 \, L_{\text{smooth}}
\label{eq:loss}
\end{equation}
where $L_{\text{vuln}}$ is MSE between simulated and observed profiles, $L_{\text{rank}} = 1 - \rho$ penalizes rank disagreement, $L_{\text{rdm}}$ encourages pairwise distance coherence, and $L_{\text{smooth}}$ regularizes adjacent-color shifts.
Effective DOF are 1--2, well below data dimensionality of 8; parameters are optimized over grids.
A candidate mechanism is selected by two criteria, not by in-sample significance: (i) \emph{structural adequacy}---the fitted parameters lie in the grid interior rather than saturating at a boundary, indicating the model has the degrees of freedom to reproduce the observed geometry; and (ii) \emph{held-out predictive loss}---using a held-out cross-validation over the HC reference pool (5 train / 2 test, $N{=}300$ resamples; and a strict 7-fold leave-one-HC-out supplement), we report the median test composite loss, and ask whether the fitted distortion predicts each held-out participant's geometry better than the no-correction baseline ($\delta\theta = 0$).
A successful LOCO precondition gate in hV4 is required for a subject to enter inference, but the LOCO criterion is a gate, not the selection metric.
HC-based false-positive calibration is reported as a descriptive diagnostic in \cref{app:hc_specificity}; it is not used as a selection rule.

%==============================================================================
\section{Results}
\label{sec:results}
%==============================================================================

\subsection{Subject-Specific Distortion Inference}

\cref{tab:main} and \cref{fig:results} present the central result: among the three candidate mechanisms, the 2-Component cortical opponent-axis model is the single model that is structurally adequate and predicts held-out HC geometry better than the no-correction baseline for both fit-recovered CVD participants, recovering a distinct sign quadrant per subject; the third participant (Sub-10, deutan) is null at both the LOCO and SRM representational levels.
The retinal R+C mechanism is rejected for both subjects by boundary saturation (\cref{sec:methods_select}, \cref{app:supporting}).
We report the mechanism class---the sign quadrant of $(\beta_s, \beta_c)$---as descriptive; the averaged loss surface for the real CVD fits is 2.1$\times$--5.5$\times$ deeper than the HC null, while absolute parameter magnitudes are not identifiable (\cref{app:supporting}).

\begin{table}[t]
  \caption{Selected distortion model per subject (2-Component cortical opponent-axis, hV4). Parameters $(\beta_s, \beta_c)$ in degrees report the mechanism class (sign quadrant), not point-resolved magnitudes. Held-out evidence: median test composite loss (5/2 HC split, $N{=}300$; lower is better) and the number of strict 7-fold leave-one-HC-out folds in which the fit beats the no-correction baseline ($\delta\theta{=}0$). The retinal R+C mechanism is rejected by boundary saturation (\cref{app:supporting}); Sub-10 is null at the LOCO and SRM levels.}
  \label{tab:main}
  \begin{center}
    \begin{small}
      \begin{tabular}{llcccc}
        \toprule
        Subject & Model & $(\beta_s, \beta_c)$ & Test loss & Folds & Mode \\
        \midrule
        Sub-08 (deutan) & 2-Comp & $(+6, -42)$ & $-2.36 \pm 2.15$ & 7/7 & ${\sim}70\%$ \\
        Sub-09 (protan) & 2-Comp & $(+2, +24)$ & $-1.54 \pm 1.42$ & 7/7 & $87.7\%$ \\
        Sub-10 (deutan) & --- & --- & --- & --- & null \\
        \bottomrule
      \end{tabular}
    \end{small}
  \end{center}
  \vskip -0.15in
\end{table}

\textbf{Sub-09 (protan).}
The 2-Component model fits Sub-09 with a dominant confusion-axis rotation and a small S-cone component ($\beta_s \approx +2^\circ$); the in-sample fit is bin-level deterministic (mode share $87.7\%$, $263/300$ resamples; held-out median test loss $-1.54$) and the fitted correction beats the no-correction baseline on all 7 leave-one-HC-out folds on the neural ($\Delta$RDM) criterion ($\Delta L = -0.472$).
We report this as a detection-and-fit result only: unlike Sub-08, Sub-09's confusion-axis \emph{sign} is not identified---the held-out per-fold oracle $\beta_c$ spans both signs (most folds negative, opposite to the in-sample $+24^\circ$), and the inferred mechanism class itself changes with the RDM metric (\cref{app:supporting}).
We therefore do not assign Sub-09 a resolved mechanism class, and Sub-09 does not anchor the divergent-correction analysis below.
The behavioral term additionally does not generalize for Sub-09 ($\Delta L = -0.55$, $4/7$ folds $\approx$ null), so the fit rests on the neural geometry alone and a direct behavioral session is the first validation priority (\cref{sec:discussion}).
The retinal R+C mechanism saturates at the grid boundary for this subject and is rejected.

\textbf{Sub-08 (deutan).}
The 2-Component model places Sub-08 in the opposite confusion-axis quadrant ($\beta_s > 0$, $\beta_c < 0$, consistent with the deutan confusion axis).
The fit has a low held-out median test loss ($-2.36 \pm 2.15$, the lowest among the structurally adequate 2-Component cells), beats the no-correction baseline on all 7 leave-one-HC-out folds for the neural criterion ($\Delta L = -0.406$) and on $5/7$ folds for the behavioral criterion ($\Delta L = -13.8$), so that neural and behavioral terms independently support the same deutan direction.
As for Sub-09, the retinal R+C mechanism saturates at the grid boundary and is rejected.

\textbf{Sub-10 (deutan).}
Sub-10 is null at both the LOCO and SRM representational levels, bounding the sensitivity of the representational levels used here; this Ishihara-confirmed case is reserved as a future-work probe of alternative representational dimensions.

\begin{figure}[t]
  \vskip 0.1in
  \centering
  \includegraphics[width=\columnwidth]{figures/fig2_loco_profiles.pdf} % TODO regenerate as RDM fig
  \caption{Subject-specific 2-Component fits (hV4). \textbf{Top}: observed, predicted, and residual $\Delta$RDM ($8{\times}8$ pairwise distance deviation from the HC mean) for each fit-recovered participant, showing that the cortical opponent-axis model reproduces the observed pairwise geometry. \textbf{Bottom}: composite loss landscape over the $(\beta_s, \beta_c)$ grid with the argmin and its bin-level stability across HC resamples. \textbf{Sub-08} (deutan): $(\beta_s, \beta_c){=}(+6^\circ, -42^\circ)$, beats no-correction on $7/7$ leave-one-HC-out folds. \textbf{Sub-09} (protan): $(\beta_s, \beta_c){=}(+2^\circ, +24^\circ)$, mode share $87.7\%$, $7/7$ folds. Parameters report the mechanism class (sign quadrant), not point-resolved magnitudes (\cref{app:supporting}). See \cref{fig:pipeline} for pipeline overview.}
  \label{fig:results}
  \vskip -0.1in
\end{figure}

\subsection{Convergent Detection, Divergent Correction}

The candidate mechanisms \emph{agree on detection}---all three identify the same two individuals as distorted---but \emph{diverge on correction}.
We make this concrete for Sub-08, whose confusion-axis sign is identified (leave-one-HC-out $\beta_c \in [-46^\circ,-38^\circ]$, no crossing of zero), so the comparison does not hinge on a non-identified parameter.
Both candidates yield \emph{feasible} corrections---the retinal Machado mechanism and the cortical 2-Component model each admit exact pre-image solutions for all eight colors---yet the corrections they prescribe point in nearly opposite directions: the per-color correction angles have cosine ${\approx}-0.4$ (range $-0.37$ to $-0.54$ across the plausible severity range $\Delta\lambda \in [2,13.5]$~nm, jackknife all-negative) and agree in sign on only $1$ of $8$ colors.
This divergence is a falsifiable prediction: the retinal account prescribes spectral pre-distortion, while the cortical account prescribes opponent-axis rescaling; because both fit the observed detection profile \emph{and} both yield feasible corrections, neither detection nor feasibility can adjudicate them---only behavioral validation can, while localizing the dominant processing level.


%==============================================================================
\section{Discussion}
\label{sec:discussion}
%==============================================================================

\paragraph{From structured distortion to intervention design.}
The LOCO vulnerability profile compresses high-dimensional voxel data into an 8-element geometry-sensitive summary.
Combined with a composite loss that decomposes geometric fidelity into per-color and pairwise components (\cref{eq:loss}), this enables distortion inference that is both interpretable and invertible into candidate corrections.
Future work could replace grid search with amortized inference (e.g., conditional normalizing flows) for real-time distortion estimation.
This approach may generalize to other conditions where sensory representations are distorted rather than absent.

\paragraph{Limitations.}
We have \emph{not} shown that corrections improve perceptual discrimination---behavioral validation is the framework's primary next step.
Three CVD participants are insufficient for population-level claims; this work is a proof-of-concept, and the fMRI requirement limits scalability.
Grid optimization on 8 data points yields elevated HC false-positive rates (15/21 model--subject pairs significant; \cref{app:hc_specificity}), which is why HC specificity is reported only descriptively and is not used as a selection rule; the pipeline targets individual profiling of \emph{known} CVD cases, not screening.
What the fit-recovered CVD cases support is mechanism-class consistency (the sign quadrant of $(\beta_s, \beta_c)$) together with cross-criterion convergence; the absolute $(\beta_s, \beta_c)$ values are \emph{not} identifiable, given a per-axis noise floor of ${\sim}20^\circ$/$25^\circ$ and a per-realization null in which 0/3 candidates pass both null sources, so the parameters are to be read as a low-dimensional descriptive embedding rather than as physiological cortical-distortion magnitudes.
Finally, behavioral generalization is asymmetric across subjects: for sub-08 both neural and behavioral terms support the deutan direction, whereas for sub-09 the fit rests on the neural geometry alone and the behavioral term does not generalize---a direct behavioral session for sub-09 is the first priority of the planned validation.


%==============================================================================
\section{Conclusion}
\label{sec:conclusion}
%==============================================================================

We present a framework that quantifies structured distortions in fMRI hue representations and fits parsimonious low-dimensional parameters, recovering subject-specific distortions that converge on detection but diverge on correction---generating falsifiable behavioral predictions as the immediate next step.

%==============================================================================
\section*{Impact Statement}
%==============================================================================

The broader implication of this work is methodological: many health conditions leave structured rather than absent signatures in high-dimensional measurements.
Mapping such signatures onto low-dimensional, invertible geometric features could turn descriptive phenotypes into actionable, individualized intervention targets.

%==============================================================================
\bibliography{SD4H_draft}
\bibliographystyle{icml2026}

%==============================================================================
%==============================================================================
\newpage
\appendix
\onecolumn

\section{Supporting Analyses}
\label{app:supporting}

\subsection{Notation}
\label{app:notation}

\begin{table}[h]
  \caption{Key notation.}
  \label{tab:notation}
  \begin{center}
  \begin{small}
  \begin{tabular}{@{}ll@{}}
    \toprule
    Symbol & Description \\
    \midrule
    $\theta(c)$, $\theta'(c)$ & Stimulus / distorted hue angle for color $c$ \\
    $\mathbf{v} \in \mathbb{R}^8$ & LOCO vulnerability profile (per-color interp.\ error) \\
    $W$ & Encoding weight matrix ($K$ channels $\to$ voxels) \\
    $\Delta\lambda$ & Retinal cone-shift magnitude (nm) \\
    $g$ & R+C cortical compensation factor ($1$: none; $2$: full; $>2$: over) \\
    $\beta_s$, $\beta_c$ & S-cone / confusion-axis angular dilation ($^\circ$) \\
    $\Delta$RDM & Pairwise distance deviation from HC mean \\
    $\rho$ & Spearman corr.\ (simulated vs.\ observed $\mathbf{v}$) \\
    \bottomrule
  \end{tabular}
  \end{small}
  \end{center}
\end{table}

\subsection{HC Specificity Calibration}
\label{app:hc_specificity}

To empirically calibrate the false-positive rate of the grid-search fitting procedure, we applied the full hV4 LOCO pipeline to all 7 HC participants, treating each as if they were a CVD candidate (best-family selected per model).

\paragraph{Results.}
\cref{tab:hc_fpr} reports the number of HC participants reaching $p < 0.05$ for each model.
The observed false-positive rate substantially exceeds the nominal 5\%: 15/21 model--subject tests are significant, compared to the expected 1.05 under the null.
The 2-component model is worst (7/7), Machado is least inflated (3/7).

\begin{table}[h]
  \caption{HC false-positive rates (hV4 LOCO, $p < 0.05$).}
  \label{tab:hc_fpr}
  \begin{center}
    \begin{small}
      \begin{tabular}{lccc}
        \toprule
        Model & DOF & HC significant & FPR \\
        \midrule
        Machado & 1 & 3/7 & 43\% \\
        R+C & 2 & 5/7 & 71\% \\
        2-Component & 2 & 7/7 & 100\% \\
        \midrule
        Total & & 15/21 & 71\% \\
        \bottomrule
      \end{tabular}
    \end{small}
  \end{center}
\end{table}

\begin{figure}[h]
  \centering
  \includegraphics[width=\textwidth]{figures/fig_hc_specificity.pdf}
  \caption{HC specificity calibration. \textbf{(a)} Per-subject p-values for each model (log scale). The red dashed line marks $\alpha = 0.05$; shaded region indicates the false-positive zone. \textbf{(b)} Aggregate false-positive rates. Dashed line: nominal $\alpha = 5\%$; dotted line: chance rate for $n{=}7$ (14.3\%). Higher model complexity (DOF) yields higher false-positive rates.}
  \label{fig:hc_specificity}
\end{figure}

\paragraph{Mechanism.}
The elevated FPR arises from two reinforcing factors:
\begin{enumerate}
\item \textbf{Grid-search flexibility on 8 data points.} With 1--2 DOF optimized over a fine grid (41--1,326 candidate parameter configurations) and rank-based loss ($L_{\text{rank}} = 1 - \rho$), most individuals---healthy or not---will have some configuration that achieves a nominally significant rank correlation with 8 elements. This is a known property of permutation tests under flexible model search.
\item \textbf{Regression to the mean.} HC individuals with low baseline LOCO performance (i.e., already poorly interpolated before any model fitting) have the most room for apparent improvement. We observe a strong negative correlation ($r = -0.894$) between HC baseline $\rho$ and the improvement $\Delta\rho$ achieved by fitting, confirming that the signal captured in HC fits is predominantly regression to the mean rather than color-specific structure.
\end{enumerate}

\paragraph{Role of this calibration.}
This false-positive calibration is reported as an honest descriptive diagnostic; it is \emph{not} a selection rule, and we do not claim specificity against healthy color vision from the present case series.
The pipeline is designed for \emph{individual profiling} of known CVD cases, not population screening.
We emphasize that the elevated HC false-positive rate is precisely why selection rests on held-out predictive loss and structural adequacy (\cref{sec:methods_select}) rather than on any in-sample significance threshold.
What can be stated descriptively about the fit-recovered CVD cases is that the 2-Component model's axes---S-cone and type-specific confusion---are grounded in known CVD physiology \citep{emery2021}, and that the fitted sign quadrant is consistent across criteria for each subject.
We make no claim that the absolute $(\beta_s, \beta_c)$ magnitudes are physiologically interpretable; per-realization null testing places 0/3 candidates past both null sources, and a $(0,0)$ algorithm validation reveals a per-axis noise floor of ${\sim}20^\circ$/$25^\circ$, so only the mechanism class is reportable.
Sub-10 is Ishihara-confirmed CVD but is null at both the LOCO and SRM representational levels used here; probing this case at alternative representational dimensions is an explicit future-work direction.

\subsection{Structural Adequacy of Candidate Mechanisms}

The retinal R+C mechanism is rejected for both CVD participants by boundary saturation: its single compensation factor $g$ collapses to the grid ceiling (sub-08 boundary rate $100\%$, $g{\to}3.0$; sub-09 boundary rate $41\%$), indicating that the model lacks a confusion-axis degree of freedom and cannot reproduce the observed geometry from the grid interior (parameter clustering at grid bounds is a recognized model-inadequacy signal).
The 2-Component model, by contrast, fits in the grid interior for both subjects, with the fitted sign quadrant consistent across the $\Delta$RDM atom computed in multiple ROIs (V1--hV4) for each participant (mechanism-class consistency; \cref{tab:full}).
These are descriptive structural-adequacy observations, not significance claims.

\subsection{Divergent Corrections at Equal Feasibility (Sub-08)}

\begin{figure}[h]
  \centering
  \includegraphics[width=\textwidth]{figures/fig_divergence_sub08.pdf} % TODO regenerate: per-color correction angles, retinal Machado vs cortical 2-Component, Sub-08 (8 hues); annotate cosine and sign-agreement
  \caption{Both candidate mechanisms yield feasible corrections, but in opposite directions (Sub-08, deutan). For each of the 8 hues, the per-color correction angle prescribed by the retinal Machado mechanism and by the cortical 2-Component model ($\beta_s,\beta_c{=}{+}6^\circ,{-}42^\circ$). Each model independently admits exact pre-image solutions for all 8 colors, yet the two correction vectors point in nearly opposite directions (cosine ${\approx}-0.4$; sign-agreement $1/8$). Feasibility therefore does not distinguish the two accounts; only their direction does, and only behavioral validation can adjudicate it.}
  \label{fig:feasibility}
\end{figure}

Both the retinal Machado mechanism and the cortical 2-Component model admit exact pre-image solutions for all 8 colors (residual $< 0.001^\circ$); feasibility alone does not separate them.
What separates them is the \emph{direction} of the prescribed correction.
For Sub-08---whose confusion-axis sign is identified (leave-one-HC-out $\beta_c \in [-46^\circ,-38^\circ]$)---the per-color correction angles of the two models have cosine ${\approx}-0.4$ (range $-0.37$ to $-0.54$ across the plausible severity range $\Delta\lambda \in [2,13.5]$~nm; leave-one-hue-out jackknife all-negative) and agree in sign on only $1$ of $8$ colors.
We anchor this comparison on Sub-08 rather than Sub-09 because Sub-09's confusion-axis sign is not identified, so its retinal-versus-cortical cosine flips sign with the (unresolved) sign of $\beta_c$ and cannot support a divergence claim.
Which correction actually improves perception requires behavioral validation.

\subsection{Model Comparison: Full Table}

\begin{table}[h]
  \caption{Candidate-mechanism comparison by structural adequacy and held-out evidence. \emph{Adequate?} marks whether the fit lies in the grid interior (R+C saturates the $g$ ceiling, so it is structurally inadequate and rejected). Held-out evidence (2-Component only): median test composite loss and the number of strict 7-fold leave-one-HC-out folds beating the no-correction baseline ($\delta\theta{=}0$) on the $\Delta$RDM criterion. R+C and Machado are listed as the rejected/dominated alternatives; no significance ($\rho$, $p$) values enter selection.}
  \label{tab:full}
  \begin{center}
    \begin{small}
      \begin{tabular}{llccccc}
        \toprule
        Subject & Model & DOF & $(\beta_s, \beta_c)$ / $g$ & Adequate? & Test loss & RDM folds \\
        \midrule
        \multirow{3}{*}{Sub-08} & Machado & 1 & $\Delta\lambda$ & dominated & --- & --- \\
        & R+C & 1 & $g{\to}3.0$ (bdy $100\%$) & \textbf{no} & --- & --- \\
        & 2-Comp & 2 & $(+6, -42)$ & \textbf{yes} & $-2.36 \pm 2.15$ & \textbf{7/7} \\
        \midrule
        \multirow{3}{*}{Sub-09} & Machado & 1 & $\Delta\lambda$ & dominated & --- & --- \\
        & R+C & 1 & $g{\approx}2.95$ (bdy $41\%$) & \textbf{no} & --- & --- \\
        & 2-Comp & 2 & $(+2, +24)$ & \textbf{yes} & $-1.54 \pm 1.42$ & \textbf{7/7} \\
        \midrule
        Sub-10 & --- & --- & --- & null & --- & --- \\
        \bottomrule
        \multicolumn{7}{p{0.95\columnwidth}}{\small R+C is rejected by boundary saturation ($g$ pinned at the grid ceiling), lacking a confusion-axis DOF. Sub-10 is null at the LOCO and SRM levels.} \\
      \end{tabular}
    \end{small}
  \end{center}
\end{table}

\subsection{Correction Feasibility Details}

\paragraph{Sub-08 (2-Component pre-image).}
All 8 colors admit exact pre-image solutions (mean $|\delta| = 46.3^\circ$; residual $< 0.001^\circ$).
The resulting 8-point lookup table is directly implementable on any digital display.

\paragraph{Sub-09 (2-Component pre-image).}
All 8 colors admit exact pre-image solutions (mean $|\delta| = 20.1^\circ$; residual $< 0.001^\circ$).

\paragraph{Retinal Machado (rejected alternative).}
The retinal Machado mechanism is rejected on structural-adequacy grounds (\cref{tab:full}), not on feasibility: evaluated as a continuous pre-image, it too admits exact solutions for all 8 colors across the plausible severity range.
Its correction is therefore feasible but mis-directed relative to the cortical fit (\cref{fig:feasibility}); a per-color arc ``collapse'' appears only under a discrete 8-bin categorical relabelling and does not hold for the deployed continuous filter, so we do not invoke it.

\end{document}
